\begin{trapbox}
	\begin{description}[style=nextline, leftmargin=2.8em, labelsep=0.5em, itemsep=0.35em]
		\item[\textbf{\textcolor{CTrap}{易错点一（退化条件）}}] 忽略$F\notin\ell$的条件，误以为当$F$在$\ell$上时$PF=PH$依然表示抛物线。
		\item[\textbf{\textcolor{CTrap}{正确理解（条件要求）：}}] 抛物线定义要求$F\notin\ell$；若$F\in\ell$，满足$PF=PH$的点$P$的轨迹是过$F$且垂直于$\ell$的直线。
		\item[\textbf{\textcolor{CTrap}{错因分析（忽略限定）：}}] 只记住了距离等量关系，忽略了$F\notin\ell$这一前提，对圆锥曲线定义的条件理解不全面。
	\end{description}

	\medskip
	\begin{description}[style=nextline, leftmargin=2.8em, labelsep=0.5em, itemsep=0.35em]
		\item[\textbf{\textcolor{CTrap}{易错点二（焦点混淆）}}] 对于标准方程$y^2=2px$，误认为焦点坐标为$(p,0)$或准线为$x=p$；对于$x^2=2py$，误认焦点为$(0,p)$。
		\item[\textbf{\textcolor{CTrap}{正确理解（标准系数）：}}] 在$y^2=2px(p>0)$中，焦点为$(\frac{p}{2},0)$，准线$x=-\frac{p}{2}$；在$x^2=2py(p>0)$中，焦点为$(0,\frac{p}{2})$，准线$y=-\frac{p}{2}$。$p$是一次项系数的一半的绝对值。
		\item[\textbf{\textcolor{CTrap}{错因分析（记忆错误）：}}] 将标准方程中$2p$与$p$的几何意义混淆，记忆抛物线标准方程时未对应好焦点位置。
	\end{description}

	\medskip
	\begin{description}[style=nextline, leftmargin=2.8em, labelsep=0.5em, itemsep=0.35em]
		\item[\textbf{\textcolor{CTrap}{易错点三（距离符号）}}] 建立距离等式时，点到直线$\ell:x=-\frac{p}{2}$的距离误写为$|x-\frac{p}{2}|$，导致方程错误。
		\item[\textbf{\textcolor{CTrap}{正确理解（直线距离）：}}] 直线$\ell:x=-\frac{p}{2}$的一般形式为$x+\frac{p}{2}=0$，点$P(x,y)$到$\ell$的距离为$|x+\frac{p}{2}|$。
		\item[\textbf{\textcolor{CTrap}{错因分析（公式记错）：}}] 记错点到直线距离公式的代入符号；对于竖直线，距离就是横坐标差之绝对值。
	\end{description}
\end{trapbox}
